\DTMsavedate{mydate}{2022-06-30}
\maketitle{Engineering Mathematics \Romannum{4}}{Applied Mathematics}{\DTMusedate{mydate}}

% ----------------------------------------------------- %
% COURSE INFO
\subsection*{Course Information}\label{course:sma3116}
\vspace{0ex}%
\noindent\parbox{0.5\textwidth}{%
    \noindent\begin{tabular}{@{}ll}
                 \textsf{Course:}          & 	Engineering Mathematics \Romannum{4}   \\
                 \textsf{Course Code:}         & 	SMA 3116    \\
                 \textsf{Credit Hours:}    & 	10
    \end{tabular}}
\vspace{2ex}\hrule\vspace{2ex}

% ----------------------------------------------------- %
% INSTRUCTOR INFO
\subsection*{Lecturer Information}
\noindent\parbox{0.5\textwidth}{%
    \noindent\begin{tabular}{@{}ll}
                 \textsf{Name:}         &    Hausitoe Nare          \\
% \textsf{Phone:}        &    \phone          \\
\textsf{Email:}        &    \href{mailto:hausitoe.nare@nust.ac.zw}{hausitoe.nare@nust.ac.zw} \\
\textsf{Qualifications:}        &  (Hon) BSc in Applied Mathematics, MSc in Operations Research
    \end{tabular}}
\vspace{2ex}\hrule\vspace{2ex}


% ----------------------------------------------------- %
% COURSE SYNOPSISDESCRIPTION
\subsection*{Course Synopsis}
This course introduces engineering students to partial differential equations, their solution using the method of separation of variables, to the power series solution of ordinary differential equations and to special functions, and to numerical methods for equation solution, curve fitting, differentiation, integration and the solution of ordinary differential equations.
\vspace{2ex} \hrule\vspace{2ex}

% ----------------------------------------------------- %
% LEARNING OBJECTIVES
\subsection*{Learning Objectives}
On completion of the course students should be able to:
\begin{outline}
    \1      Solve simple partial differential equations by direct integration or by the use of the techniques of ordinary differential equations.

    \1  	Solve eigenvalue and eigenfunction problems that arise in simple boundary value problems.

    \1  	Use the method of separation of variables to solve boundary value problems with the Laplace, Diffusion and Wave equations.

    \1      Use the power series method to solve ordinary differential equations in the neighbourhood of an ordinary point.

    \1      Identify and classify the singular points of an ordinary differential equation.

    \1      Use the method of Frobenius to solve ordinary differential equations in the neighbourhood of an regular singular point.

    \1      Solve Bessel’s equation and others to identify a solution as a particular special function.

    \1      Understand the concept of different types of error encountered in numerical methods.

    \1      Understand and use a variety of methods to fit a curve to a set of points.

    \1      Understand and use a variety of methods, including the Three Point Method, to differentiate a function of one variable.

    \1      nderstand and use a variety of methods, including Simpson’s rule, to evaluate definite integrals.

\end{outline}
\vspace{2ex}\hrule\vspace{2ex}


\subsection*{Learning Requirements}
Students are expected to have completed and passed the following courses:
\begin{outline}
    \1      SMA 1116

    \1      SMA 1216

    \1      SMA 2116

    \1      SMA 2217


\end{outline}
\vspace{2ex}\hrule\vspace{2ex}



% ----------------------------------------------------- %
% COURSE TOPICS
\subsection*{Topics}
\begin{outline}
    \1 Ordinary Differential Equations
    \begin{itemize}
        \item Power series solutions
        \item Frobenius method
        \item Special functions and their properties
        \item Legendre polynomials
        \item Bessel functions
    \end{itemize}
    \1 Partial Differential Equations
    \begin{itemize}
        \item Solutions of partial differential equations
        \item Method of separation of variables
    \end{itemize}
    \1 Numerical Methods
    \begin{itemize}
        \item Errors, relative and absolute
        \item The solution of non-linear equations
        \item The solution of linear systems
        \item Interpolation and polynomial approximation
        \item Curve fitting
        \item Numerical differential and integration
        \item Approximate solution of differential equations
    \end{itemize}
\end{outline}
\vspace{2ex} \hrule\vspace{2ex}

% ----------------------------------------------------- %
% Students Assenssment
\subsection*{Assessment of Students}
\begin{outline}
    \1 Students are assessed through:
    \begin{enumerate}
        \item course work consisting of two in-class tests, test on one Ordinary Differential Equations and Partial Differential Equations and test two on Numerical Methods.
        \item end of semester final examination with Section A of 40 marks and Section B of 60 marks.
    \end{enumerate}

\end{outline}
\vspace{2ex}\hrule\vspace{2ex}

% ----------------------------------------------------- %
% Grade Evaluation
\subsection*{Course Evaluation and Grading Policies}
Grades are based on:
Continuous assessment (\qty{25}{\percent}),
Final Exam (\qty{75}{\percent})
\vspace{2ex}\hrule\vspace{2ex}

% ----------------------------------------------------- %
% EQUIPMENT
% https://sites.udel.edu/computing-purchases/personal-specs/
% \subsection*{Essential Equipment}

% \vspace{2ex}\hrule\vspace{2ex}

% Policies and Other information
\subsection*{Policies and Other Information}
% Conduct
\subsection*{Key Text}
\begin{itemize}
    \item Elementary Differential Equations and Boundary Value Problems, Eighth edition by  W.E. Boyce and  R.C. DiPrima
    \item Numerical Analysis, 7th Edition by R.L. Burden and Douglas L. Faires
    \item Partial Differential Equations, 2nd Edition by E. DiBenedetto
    \item Elements of Partial Differential Equations by P. Drabek and G. Holubova
    \item Elementary Differential Equations with Boundary Value Problems by H. Edwards, and D. Penney
    \item Beginning Partial Differential Equations by P.V. O’Neil
    \item Ordinary Differential Equations by M. Tennenbaum and H. Pollard
\end{itemize}
\vspace{2ex}\hrule
\vspace{2ex}\hrule\vspace{2ex}